\reportpage{\dataheading}{
\subheading{本节原始输出证据}
以下为本次抓包的TShark字段、原始字节或命令日志摘录；完整文本与对应pcapng随附件保存。
\ev{fig:ping-request}{ping-request}{普通Ping请求的ICMP、IPv4字段与原始字节（帧1）（原始输出摘录）}
}

\reportpage{\datacontinued}{
\subheading{本节原始输出证据}
以下为本次抓包的TShark字段、原始字节或命令日志摘录；完整文本与对应pcapng随附件保存。
\ev{fig:ping-reply}{ping-reply}{对应回显应答的字段与数据（帧2）（原始输出摘录）}
}

\reportpage{\datacontinued}{
\subheading{二、普通Ping的ICMP报文（题1、2）}
本节普通测试均引用\capfile{01-ping-default}，本机为\code{10.135.133.108}，目标为\code{202.115.32.43}。该文件共10帧，五组请求与应答，命令报告5发5收、丢包率0\%。

\qa{1. 本次Ping使用哪些ICMP报文？Type和Code是什么？}{
\figref{fig:ping-request}中的请求为Echo Request，\code{Type=8, Code=0}；\figref{fig:ping-reply}中的应答为Echo Reply，\code{Type=0, Code=0}。帧1从本机发往目标，帧2端点反向；其余四组同样配对。普通测试中没有捕获到差错报告。
}
\qa{2. 请求数据有多少字节？内容及对应应答有何变化？}{
帧1的ICMP数据为\textbf{56字节}，另有8字节ICMP首部，故ICMP报文共64字节，外层IPv4总长84，Ethernet帧长98。数据前8字节为\code{6a eb 0c 64 00 02 ab dc}，被Wireshark解释为发送时间戳；其后48字节按\code{08,09,0a,\ldots,36,37}递增。低位字节含控制字符，例如\code{09}为制表、\code{0a}为换行；\code{20}为空格，\code{30--37}为字符0至7，不能将整段称作普通可读文本。

\figref{fig:ping-request}显示的\code{Data (48 bytes)}是时间戳以外的剩余部分，\textbf{不是完整ICMP数据仅48字节}。帧2逐字节返回同一56字节数据，含相同时间戳；变化的是Type、ICMP校验和及外层IPv4首部等字段。其余四组也均通过数据字节相等验证。
}
}

\reportpage{\datacontinued}{
\subheading{本节原始输出证据}
以下为本次抓包的TShark字段、原始字节或命令日志摘录；完整文本与对应pcapng随附件保存。
\ev{fig:ping-time}{ping-time}{普通Ping命令输出与首次应答时间记录（原始输出摘录）}
}

\reportpage{\datacontinued}{
\subheading{三、往返时间、校验和与报文配对（题3--5）}
\qa{3. 首次收到应答的请求，其RTT是多少？与Ping如何比较？}{
\capfile{01-ping-default}首次成功配对为帧1和帧2，见\figref{fig:ping-request}、\figref{fig:ping-reply}。捕获时间分别为北京时间\code{18:43:48.175240}和\code{18:43:48.180669}，相减得
\[
 (48.180669-48.175240)\times1000=5.429\ \mathrm{ms}.
\]
\figref{fig:ping-time}中相应\code{icmp\_seq=0}的命令结果为5.759毫秒，相差0.330毫秒。抓包时间是在捕获接口处记录，Ping从应用/协议栈发送、接收处理计时，边界和调度不同，二者不要求完全相等。不能把数据中的发送时间戳与捕获时间混为一类。

\par\noindent\begin{tabularx}{\linewidth}{|c|c|c|X|X|}\hline
Seq & 请求帧 & 应答帧 & 抓包RTT/ms & Ping/ms\\\hline
0&1&2&5.429&5.759\\\hline
1&3&4&5.159&5.572\\\hline
2&5&6&4.702&5.182\\\hline
3&7&8&4.801&5.036\\\hline
4&9&10&4.628&5.086\\\hline
\end{tabularx}\par
}
\qa{4. IPv4已有Header Checksum，为什么ICMP仍需Checksum？}{
IPv4校验和只覆盖\textbf{IPv4首部}，不检查其载荷；ICMP校验和覆盖\textbf{整个ICMP首部与ICMP数据}，不含外层IPv4首部，也不使用UDP式伪首部。因此二者保护的范围不同。\figref{fig:ping-request}帧1的IPv4首部20字节，Header Checksum为\code{0x8bc2}；ICMP覆盖64字节，Checksum为\code{0xe8da}。\figref{fig:ping-reply}帧2分别为\code{0x0cc9}和\code{0xf0da}。独立按16位反码和复算，两种校验均正确；GUI的IPv4“Unverified”仅表示默认未开启验证。
}
\qa{5. Identifier和Sequence Number的作用及本次对应关系？}{
Identifier用于区分回显会话，Sequence Number用于区分该会话的各次请求；应答复制这两个字段。\figref{fig:ping-request}、\figref{fig:ping-reply}采用网络序BE值，Identifier均为\code{0x00f4}=244；五次Seq依次为0、1、2、3、4，每组请求与应答相同。结合反向IP端点、相同数据和先后时间共同配对。IPv4 Identification如帧1的\code{0x7455}是\textbf{另一字段}，不能当作ICMP Identifier。
}
}

\reportpage{\datacontinued}{
\subheading{本节原始输出证据}
以下为本次抓包的TShark字段、原始字节或命令日志摘录；完整文本与对应pcapng随附件保存。
\ev{fig:large-first}{large-first}{2500字节Ping首片IPv4字段（帧1）（原始输出摘录）}
\ev{fig:large-reassembly}{large-reassembly}{第二片偏移与2508字节ICMP重组结果（原始输出摘录）}
}

\reportpage{\datacontinued}{
\subheading{四、2500字节数据与IPv4分片（题6）}
\qa{6. 所选请求是否分片？各片长度、偏移、MF及重组长度？}{
命令使用macOS等价参数\code{-s 2500}。\capfile{02-ping-2500}共20帧，5个请求和5个应答均各分为两片。\figref{fig:large-first}、\figref{fig:large-reassembly}所选首个请求由帧1、2构成：源/目的地址、Protocol=1及IPv4 Identification \code{0xfb81}相同；MF和偏移连续，支持将其识别为同一个IPv4数据报的分片。

\par\noindent\begin{tabularx}{\linewidth}{|c|X|X|c|c|c|}\hline
帧 & IPv4总长 & IPv4负载 & MF & 偏移字段 & 字节偏移\\\hline
1&1500&1480&1&0&0\\\hline
2&1048&1028&0&185&1480\\\hline
\end{tabularx}\par
各片IPv4首部均20字节、DF=0。偏移字段以8字节为单位，第二片\code{185}\,$\times8=1480$字节；该界面显示185时不能直接读作185字节。首片负载1480由8字节ICMP首部和1472字节数据组成；第二片的1028字节全为后续数据，\textbf{没有重复ICMP首部}。

重组后的IPv4负载为$1480+1028=2508$字节，即8字节ICMP首部加2500字节ICMP数据；若未分片，IPv4总长为2528字节。线上两片总长之和2548多出一次20字节IPv4首部，不能称为ICMP数据长度。首个应答由帧3、4重组，IPv4 ID为\code{0xf849}，长度及偏移结构相同；它与请求的2500字节数据逐字节一致。10个重组ICMP报文均通过独立校验和验证。

首个重组请求/应答的ICMP Identifier为\code{0x1cf4}，Seq=0，Checksum分别为\code{0x2b3b}/\code{0x333b}。Wireshark完成重组后在帧2、4解释ICMP，故其RTT为$18{:}43{:}53.217385-18{:}43{:}53.211824=5.561$毫秒；采用两首片时刻会得5.590毫秒，计时边界应注明。
}
}

\reportpage{\datacontinued}{
\subheading{本节原始输出证据}
以下为本次抓包的TShark字段、原始字节或命令日志摘录；完整文本与对应pcapng随附件保存。
\ev{fig:ttl3}{ttl3}{TTL=3的Time Exceeded与原报文引用（帧2）（原始输出摘录）}
}

\reportpage{\datacontinued}{
\subheading{五、TTL耗尽与差错报文（题7）}
\qa{7. TTL=2及补测TTL=1收到什么响应？由谁发送？}{
\figref{fig:ttl-commands}保留各轮命令结果。\capfile{03-ping-ttl02}和\capfile{04-ping-ttl01}均只有5个本机请求，没有捕获目标回显应答，也没有捕获Time Exceeded；两轮Ping均报告0收、100\%丢包。因此这两轮没有可列出的“实际响应Type/Code及发送者”。不能将请求超时直接认定为目标宕机，路由设备可能限速或不返回ICMP差错。

继续逐步增加TTL后，\capfile{05-ping-ttl03}共10帧，5个请求与5个Time Exceeded。\figref{fig:ttl3}所选帧1初始TTL=3，目标仍为\code{202.115.32.43}；帧2由中间设备\code{202.115.39.102}发给本机\code{10.135.133.108}，外层ICMP为\code{Type=11, Code=0}，表示报文在转发过程中TTL耗尽，并非目标的Echo Reply。其外层IPv4 TTL为253，是\textbf{差错报文自己的TTL}。

帧2外层IPv4总长56字节，包括20字节IPv4首部和36字节ICMP；ICMP由8字节差错首部及28字节原报文引用构成：\textbf{原IPv4首部20字节+原ICMP首部8字节}。Unused为全零，外层Checksum为\code{0xac85}，复算正确。

引用内层仍为本机到目标，IPv4总长字段84、Identification \code{0xe0fe}、Protocol=1，所引IPv4 TTL实际为1；这不是外层TTL，也不能擅自改成0。内层Echo字段为\code{Type=8, Code=0}、Checksum \code{0x2f85}、Identifier \code{0x10f5}、Seq=0，与帧1一致。引用仅含原ICMP的前8字节，原56字节数据未返回，故不能仅据此重新计算内层完整ICMP校验和。

该对差错反馈的捕获时差为15.868毫秒；另外四对为帧3/4、5/6、7/8、9/10，Seq=1至4，反馈分别为7.121、8.068、7.363、7.369毫秒。Ping仍显示“0 packets received”，因为这些是中间设备的\textbf{差错报告}，没有一个是目标回显应答。附件的字段导出中重复出现的IP、TTL、Type分别属于外层与引用层，分析时需分层解释。
}
}

\reportpage{\datacontinued}{
\subheading{本节原始输出证据}
以下为本次抓包的TShark字段、原始字节或命令日志摘录；完整文本与对应pcapng随附件保存。
\ev{fig:ttl-commands}{ttl-commands}{TTL=1、2、3、4各轮Ping原始命令结果摘录（原始输出摘录）}
}

\reportpage{\datacontinued}{
\subheading{六、最小成功TTL及历史攻击（题8、9）}
\qa{8. 本轮首次收到目标回显应答时的初始TTL是多少？}{
如\figref{fig:ttl-commands}所示，TTL=1、2未收到应答，TTL=3仅收到\figref{fig:ttl3}的中间设备Time Exceeded，TTL=4首次收到目标\code{202.115.32.43}的Echo Reply，故本次逐TTL测试的最小成功初始值为\textbf{4}。

\capfile{05-ping-ttl04}有5组请求与应答，帧1/2、3/4、5/6、7/8、9/10对应Seq=0至4；请求初始TTL均4，ICMP Identifier均\code{0x64f5}；应答Type=0、Code=0，发送者为目标，应答到达时TTL均61。命令5发5收，RTT最小/平均/最大为3.870/4.777/5.521毫秒。本次已经在TTL=4满足停止条件，无须继续至30。

TTL限制的是\textbf{所发报文的转发存活范围}；该成功阈值受当前路径、路由变化、负载均衡和设备ICMP策略等影响，只对应本机、本目标和本轮测量。应答所带TTL是返回报文独立设置并递减的值，不能把61或本次阈值4当作其他主机、其他时刻固定不变的距离。
}
\qa{9. 举出两类利用ICMP或Ping的历史攻击，并与实验比较。}{
\textbf{Smurf反射放大：}将受害者地址伪作源地址，对允许响应的网络广播地址发送Echo Request，诱使多个主机向受害者回显，集中消耗链路与处理资源。参见\href{https://www.rfc-editor.org/rfc/rfc4732.html\#section-3.1}{RFC 4732第3.1节}。本次\figref{fig:ping-request}为本机真实源地址到单一目标的有限单播请求，没有伪造源地址或广播放大。

\textbf{Ping of Death：}利用旧协议栈对异常超长分片重组处理的缺陷，使重组后的IPv4报文超过65535字节上限，引起拒绝服务。参见\href{https://www.rfc-editor.org/rfc/rfc4732.html\#section-2.1.1}{RFC 4732第2.1.1节}及\href{https://www.cve.org/CVERecord?id=CVE-1999-0128}{CVE-1999-0128}。本次\figref{fig:large-reassembly}仅重组为2528字节IPv4报文，分片偏移连续、数据完整；2500字节数据大于MTU引起的\textbf{合法分片本身并非攻击}。实验各轮仅发送5个请求，目的为观察协议行为，未实施上述攻击。
}
}
